\chapter{Gluon PDFs: Where the Two Lines Converge}

Gluon distributions are where the field's difficulties concentrate:
adjoint operators, larger statistical noise, single-charge-conjugation
mixing with quarks in the singlet channel --- and the largest theory
uncertainty in global fits. The five papers of this chapter supplied the
theory for the library's gluon-PDF program and fused the gradient-flow
and large-momentum lines.

% ----------------------------------------------------------------------
\section{Matching Gluon Quasi-PDFs}
\papercard{J.-H.\ Zhang, X.\ Ji, A.\ Sch\"afer, W.\ Wang, S.\ Zhao}
{Accessing gluon parton distributions in large momentum effective theory}
{Phys.\ Rev.\ Lett.\ \textbf{122}, 142001 (2019)}
{arXiv:1808.10824 | DOI: 10.1103/PhysRevLett.122.142001}
\whyselected{The theoretical engine of every gluon quasi-PDF calculation
in the library ($\sim$9/99).}
\physics{Using an auxiliary adjoint heavy-quark field coupled through the
Yang--Mills Lagrangian as a gauge-invariant regulator, the paper proves
that all power divergences of gluon quasi-PDF operators reside in the
Wilson-line self-energy and are removable by mass renormalization. It then
constructs the independent multiplicatively renormalizable correlators ---
four unpolarized, three helicity --- computes their one-loop matching to
light-cone gluon PDFs, and states the LaMET factorization formula for the
gluon channel. The classification also clarifies which operator
combinations are free of mixing with quark singlets at leading twist.}
\impact{The library's Fan et al.\ gluon quasi-PDF letter and the Lin-group
LaMET gluon series implement precisely this framework; the hybrid-renormalization
papers adapt its subtraction to finite-flow operators.}
\cardend

% ----------------------------------------------------------------------
\section{Gluon Pseudo-Distributions from Short Distances}
\papercard{I.\ Balitsky, W.\ Morris, A.\ Radyushkin}
{Gluon pseudo-distributions at short distances: forward case}
{Phys.\ Lett.\ B \textbf{808}, 135621 (2020)}
{arXiv:1910.13963 | DOI: 10.1016/j.physletb.2020.135621}
\whyselected{Theory hub of the entire HadStruc gluon line ($\sim$10/99):
every gluon pseudo-ITD extraction in this library quotes its matching
kernel.}
\physics{The paper classifies bilocal double-gluon operators and isolates
the twist-2 invariant amplitudes entering forward matrix elements; using
the background-field method and heat-kernel techniques it computes the
short-distance expansion of the gluon pseudo-ITD in an explicitly
gauge-invariant form,
\begin{equation}
  \mathcal{M}_g(\nu,z^2)\big|_{|z^2|\ \mathrm{small}}
  = C_g(\nu,z^2,\mu)\otimes g(x,\mu) + \mathcal{O}(z^2),
\end{equation}
separating ultraviolet logarithms from DGLAP evolution and supplying the
reduced-ITD matching relation. The companion paper (PRD \textbf{105},
014008 (2022), arXiv:2111.06797, also cited widely in the library)
provides complete expressions for all projections and the quark-mixing
terms.}
\impact{Khan et al., Egerer et al., Salas-Chavira et al., Fan--Lin and
the physical-point/continuum extractions all convert their matrix elements
through these kernels --- they are to gluon pseudo-PDFs what Xiong et al.\
is to the nonsinglet quark channel.}
\cardend

% ----------------------------------------------------------------------
\section{Smeared Quasi-PDFs: The Flow Meets Partons}
\papercard{C.\ Monahan, K.\ Orginos}
{Quasi parton distributions and the gradient flow}
{JHEP \textbf{03} (2017) 116}
{arXiv:1612.01584 | DOI: 10.1007/JHEP03(2017)116}
\whyselected{Original article of this library and the explicit bridge
between Chapters~3 and 4 ($\sim$5/99 within its niche, universal within
the smeared-operator subline): it imports the flow's finiteness theorem
into parton physics.}
\physics{Replace the bare Wilson-line operator of \cref{eq:quasipdf} by a
version built from flowed fields $B_\mu(t)$, $\chi(t)$ at small flow time
$t$: since flowed products are UV finite (\cref{eq:flow}), the matrix
element has a smooth continuum limit without manual subtraction of linear
divergences. With ringed fermions the wave-function renormalization stays
multiplicative; the small-$t$ expansion relates moments of the smeared
quasi-PDF to light-cone moments with computable coefficients obeying a
DGLAP-like system. The follow-up one-loop analysis by Monahan (PRD
\textbf{97}, 054507 (2018), in the library) supplies the matching kernels
and confirms that smearing leaves infrared structure untouched.}
\impact{This is the conceptual source of the library's
``Gradient Flow for PDFs'' pion application and of Shindler's flowed
twist-2 moment proposal: once operators are smeared by the flow, the
renormalization battle of Chapter~5 simplifies into ordinary matching ---
at the price of a new scale $t$ whose systematic treatment the NNLO
bilinear papers provide.}
\cardend

% ----------------------------------------------------------------------
\section{Self-Renormalization}
\papercard{Y.-K.\ Huo et al.\ (Lattice Parton Collaboration)}
{Self-renormalization of quasi-light-front correlators on the lattice}
{Nucl.\ Phys.\ B \textbf{969}, 115443 (2021)}
{arXiv:2103.02965 | DOI: 10.1016/j.nuclphysb.2021.115443}
\whyselected{Original article of this library and the workhorse scheme of
its LP3 program ($\sim$6/99): used by the transversity, DA and TMD papers.}
\physics{Exploit the structure of the divergence itself: write the
renormalized matrix element as
\begin{equation}
  Q_R(z,P^z;\mu)=Z^{-1}(z,a,m_0)\;Q_{\rm lat}(z,P^z,a),
  \qquad Z(z,a)\propto e^{\,-c\,z/a}\times(\text{logs}),
\end{equation}
and determine both $c$ and the residual log factor by fitting
multi-spacing data on $Q_{\rm lat}$ itself --- no external vertex
calculation, no Landau-gauge dependence, no auxiliary field. The paper
verifies universality of the extracted $Z$ across hadron momenta and
channels, and documents sizable nonperturbative contamination in
RI/MOM-style vertices at long distances, motivating short-distance-only
usage.}
\impact{Within the library this is the scheme behind several flagship
results (nucleon transversity in the continuum limit; pion/kaon DAs;
TMD wave functions); its diagnostics fed directly into the
hybrid-scheme design below.}
\cardend

% ----------------------------------------------------------------------
\section{Hybrid Renormalization}
\papercard{X.\ Ji, Y.-S.\ Liu, A.\ Sch\"afer, W.\ Wang, Y.-B.\ Yang,
J.-H.\ Zhang, Y.\ Zhao}
{Hybrid renormalization scheme for quasi light-front correlations in
large-momentum effective theory}
{Nucl.\ Phys.\ B \textbf{964}, 115311 (2021)}
{arXiv:2008.03886 | DOI: 10.1016/j.nuclphysb.2021.115311}
\whyselected{Original article of this library; the current standard
scheme of its LP3 program ($\sim$10/99): the pragmatic synthesis after a
decade of battle.}
\physics{Split the separation axis into regimes. For $z<z_c$ use a ratio
of matrix elements (or RI/MOM) at short distance, where perturbation
theory is trustworthy and nonperturbative pollution is absent; for
$z>z_c$, subtract the linearly divergent Wilson-line self-energy through
an auxiliary-field mass term matched continuously onto the short-distance
regime, then convert to $\overline{\rm MS}$. The scheme carries one
parameter $m_0^2$ (equivalently $z_c$) with intrinsic $\mathcal{O}(a)$
ambiguity, for which the paper prescribes determination procedures ---
fitting across spacings, or fixing via the cusp anomalous dimension.
Combined with momentum-space resummation of $\ln(p_z/\mu)$ logs
(Su et al., NPB \textbf{991}, 116201 (2023), arXiv:2209.01236, in the
library) and renormalon-consistent power corrections (Zhang et al.,
PLB \textbf{844}, 138081 (2023)), it defines the modern pipeline.}
\impact{Every late-stage LP3 result in this library --- gluon PDFs,
Boer--Mulders, heavy-meson LCDAs, TMD soft functions --- runs on this
scheme or its self-renormalization variant; the two papers together are
the most-quoted methodological pair of the corpus's second half.}
\cardend
